%%% Generated by SAPlugin v10.0 on Tue May 19 01:09:35 CEST 2026

% (optional)
\chapter{Introduction}\label{sec:introduction}

This report documents the second iteration of the \textsc{pms} architecture.
The goal of the iteration is \emph{not} to add functionality for its own
sake, but to extend the initial architecture so that it \emph{fully supports}
two additional quality attribute scenarios (QASs) without undermining the
QAS-driven decisions already taken for \texttt{P2: Risk estimation by clinical
models} (Medium) and \texttt{M2: New or optimized clinical models} (High).

\paragraph{Selected scenarios.}
From the catalogue of provided QASs the following two scenarios are
addressed:
\begin{itemize}
  \item \textbf{Av2: Communication between patient gateways and the
        \textsc{pms} backend (High).}
        The \textsc{pms} availability of the end-to-end sensor-data path
        depends on (i) the external telecom channel, (ii) the patient
        gateway device, and (iii) the internal backend subsystems that
        ingest, route, and persist the readings. The scenario requires the
        system to \emph{prevent}, \emph{detect}, and \emph{resolve}
        failures along this path under the quantitative measures laid down
        in Section~3.5 of the requirements (99.5\,\%/80\,\%/64\,Kb/s SLA,
        $T{+}5$\,min missing-data detection, 5\,s internal-failure
        detection, 2/5/10\,min stakeholder notification per risk level,
        5\,min redeployment, 24\,h local buffering on the gateway).
  \item \textbf{P1: Data exchange with physicians (Medium).}
        The throughput-bounded scenario in which the \textsc{pms} handles
        physician-initiated reads (UC6), on-demand consultations (UC9), risk
        re-configuration (UC7), risk-level updates (UC8), and
        event-driven notifications dispatched to a per-physician inbox for
        registered physicians (UC5). The
        scenario imposes per-tier latency SLAs (2/5/10\,s reads), a
        100\,notifications/min throughput target with strict priority order
        (red\,$>$\,yellow\,$>$\,green) and a non-interference requirement
        with respect to the risk-estimation pipeline already covered by P2.
\end{itemize}

\paragraph{Scope of the extension.}
The original architecture contained the placeholder component
\texttt{OtherFunctionality} that absorbed every responsibility outside of
the risk-estimation pipeline. The extension replaces this placeholder with
concrete, QAS-driven components, organised in two cohesive subsystems:
\begin{enumerate}
  \item an \emph{Av2-subsystem} that owns the patient gateway, the backend
        ingress chain (gateway $\rightarrow$ broker $\rightarrow$
        ingestion), the monitoring/recovery tier, and the administrative
        redeployment control plane;
  \item a \emph{P1-subsystem} that owns the physician-facing gateway, the
        tiered query/command services, the event-driven notification
        dispatcher with a per-physician inbox, the patient-data service
        (sensor data + EHR proxy),
        and the patient-status cache.
\end{enumerate}
Both subsystems are layered on top of the unchanged risk-estimation
pipeline (\texttt{RiskEstimationScheduler}, \texttt{RiskEstimationProcessor},
\texttt{RiskEstimationCombiner}, \texttt{ClinicalModelCache},
\texttt{ClinicalJobCreator}, \texttt{ClinicalModelDB}). The original P2/M2
decisions are preserved verbatim; this report makes their interactions with
the new subsystems explicit (see Chapter~\ref{ch:discussion}).

\paragraph{Reasoning style.}
The rationale below is written in the spirit of the Architecture Trade-off
Analysis Method (ATAM). For each architectural decision we identify the
\emph{sensitivity points} introduced or affected, the \emph{trade-offs}
with respect to other quality attributes, and the impact on previously
documented decisions. We do not perform a full ATAM analysis; we use its
vocabulary to discipline the discussion.


% ============================================================
% Part 2: QAS Decisions
% ============================================================
\chapter{QAS Decisions}
\label{ch:qas-decisions}

\section{Av2: Communication between patient gateways and \textsc{pms}
backend (H)}
\label{sec:av2}

\subsection*{Key Decisions}
\begin{description}
  \item[(AD-Av2-1)] \textbf{Edge-resilient patient gateway.}
        The \texttt{Av2PatientGateway} owns local buffering for at least
        24\,h of sensor data, an exponential-backoff retry strategy, an
        explicit \emph{degraded-mode} state machine, and a backup
        communication channel (SMS) for emergency notifications.
  \item[(AD-Av2-2)] \textbf{Single ingress with broker-decoupled
        persistence.}
        The \texttt{Av2CommunicationGateway} is the sole backend ingress
        point; it authenticates, routes, and acknowledges, and writes to
        the durable \texttt{Av2MessageBroker} from which the
        \texttt{Av2DataIngestionService} drains asynchronously into the
        \texttt{Av2SensorDataRepository}.
  \item[(AD-Av2-3)] \textbf{Heartbeat-and-health detection within Av2
        bounds.}
        Every transmission updates the \texttt{Av2HeartbeatTracker}; the
        \texttt{Av2MissingDataDetector} polls it against the
        risk-level-calibrated expected interval so that an overdue gateway
        is flagged within $T{+}5$\,min. Internal subsystems expose
        health-check endpoints polled by the
        \texttt{Av2InternalFailureDetector} at $\geq 1$\,Hz so that any
        subsystem crash is detected within 5\,s.
  \item[(AD-Av2-4)] \textbf{Risk-aware failure classification and
        notification.}
        Detected failures are enriched by the \texttt{Av2FailureClassifier}
        with the affected patient's risk level (read via
        \texttt{Av2PatientStatusOverview}), which determines an escalation
        deadline of 2/5/10\,min (HIGH/MEDIUM/LOW). The
        \texttt{Av2StakeholderNotifier} must meet the deadline in
        $\geq 90\,\%$ of cases.
  \item[(AD-Av2-5)] \textbf{Administrator-controlled recovery.}
        The \texttt{Av2RedeploymentCoordinator} (driven by the
        \texttt{Av2AdminPortal}) bounds redeployment/rollback to 5\,min;
        the \texttt{Av2RecoverySyncService} then drains the broker backlog
        and instructs affected gateways to flush their local buffer in
        chronological order.
  \item[(AD-Av2-6)] \textbf{SLA observability.}
        The \texttt{Av2SLAMonitor} aggregates per-gateway uptime, coverage
        and bandwidth samples against the SLA thresholds (99.5\,\% per
        year, 80\,\% coverage, 64\,Kb/s) and raises a
        \texttt{SLA\_VIOLATION} failure event when any rolling threshold
        is breached.
\end{description}

\subsection*{Employed tactics and patterns}
\textbf{Availability tactics:} \emph{ping/echo} (heartbeats and health
checks), \emph{heartbeat with adaptive interval}, \emph{exception
detection}, \emph{retry with exponential back-off}, \emph{redundancy via
spare channel} (SMS), \emph{state resynchronization} (post-failure broker
drain + gateway flush), \emph{rollback} (administrator-triggered
redeployment), \emph{degraded operation} (gateway degraded-mode state
machine), \emph{transactional buffer} (durable broker between ingress and
persistence).
\textbf{Patterns:} \emph{gateway}, \emph{message broker / store-and-forward
queue}, \emph{circuit-breaker variant} on the routing path,
\emph{publish/subscribe} for failure escalation,
\emph{health-check endpoint} pattern, \emph{ack/clear} pattern between
backend and edge buffer.

\subsection*{Architectural rationale}

\paragraph{(AD-Av2-1) Edge-resilient patient gateway.}
The first failure mode in the QAS is that the external telecom channel or
the patient gateway itself becomes unavailable. The architecture pushes the
first line of defence onto the gateway device because that is the only
location at which sensor data still exist when the channel is broken: any
purely backend-side mitigation would arrive too late. We therefore
decompose the \texttt{Av2PatientGateway} into the following internal
modules:
\begin{itemize}
  \item \texttt{Av2SensorCommunicator} reads raw readings via
        \texttt{Av2SensorDataAcquisition} and pushes them into the
        in-memory \texttt{DataBuffer};
  \item \texttt{Av2SensorDataTransmitter} consumes the buffer and, in
        \emph{normal} mode, calls \texttt{Av2GatewayDataIngress} on the
        backend \texttt{Av2CommunicationGateway};
  \item \texttt{Av2ServerConnectivityMonitor} continuously probes the
        backend; if the backend is unreachable the
        \texttt{Av2DegradeModeManager} transitions the gateway into
        \emph{degraded} mode and timestamps the transition;
  \item in degraded mode the transmitter hands its readings to
        \texttt{Av2LocalBufferingRepository} instead, which guarantees
        storage of \emph{at least 24\,h} of readings at the configured
        transmission rate -- directly satisfying Av2 resolution response
        measure 3 (``the patient gateway is able to store the sensor data
        readings and unsent notifications for the past 24 hours'');
  \item \texttt{Av2RetryAndSynchronizationService} reads pending entries
        and retries delivery with exponential back-off (Av2 resolution
        response 3, ``keeps retrying to deliver sensor data readings in a
        proper fashion, e.g., using exponential back-off'');
  \item upon successful delivery the \texttt{Av2AcknowledgementService}
        sends an explicit \texttt{Av2DeliveryAck}, allowing the retry
        service to clear the corresponding entries from the buffer.
        Clearing requires an acknowledgement, which is what closes the
        loop and prevents data loss on a crash mid-flush;
  \item \texttt{Av2BatteryMonitor} and
        \texttt{Av2NetworkConnectivityMonitor} drive the
        \texttt{Av2PatientAlertService}, which warns the patient when the
        battery drops to $\leq 15\,\%$ or connectivity is poor / absent --
        satisfying Av2 \emph{prevention} response measures verbatim;
  \item the \texttt{Av2EmergencyNotificationService} consults
        \texttt{Av2DegradeModeStatus} and, when in degraded mode, delegates
        to the \texttt{Av2BackupChannelManager}, which routes the emergency
        notification through SMS via the external
        \texttt{SMSGatewayProvider} and the backend
        \texttt{Av2EmergencyBackupReceiver}. This satisfies the ``uses a
        back-up communication channel for emergency notifications, e.g.,
        SMS'' response.
\end{itemize}

\paragraph{(AD-Av2-2) Single ingress with broker-decoupled persistence.}
On the backend side we deliberately reject a fan-out topology in which
patient gateways speak directly to the \texttt{Av2DataIngestionService} or
to the existing \texttt{RiskEstimationScheduler}. Instead, all
primary-channel traffic enters the system at one point -- the
\texttt{Av2CommunicationGateway} -- and is dropped into a durable
\texttt{Av2MessageBroker} from which the \texttt{Av2DataIngestionService}
drains asynchronously. Two reasons drive this:
\begin{enumerate}
  \item \emph{Failure isolation for response 6 of P1.} A transient outage
        of the ingestion service or the patient-record persistence layer
        must not propagate back to the gateways, nor must it stall the
        risk-estimation pipeline. The broker is the explicit
        \emph{transactional buffer} that decouples ingress rate from
        downstream processing rate, which is also what the P1 scenario
        requires for response 6 (``A large number of notifications to be
        sent or requests arriving to the system should have no impact on
        (a) the processing of incoming raw sensor reading data; or (b)
        ongoing risk estimations.'').
  \item \emph{Detectability.} Both the \texttt{Av2CommunicationGateway}
        and the \texttt{Av2MessageBroker} expose health-check endpoints
        (\texttt{Av2CommunicationGatewayHealthCheck},
        \texttt{Av2MessageBrokerHealthCheck}). The
        \texttt{Av2InternalFailureDetector} polls them at $\geq 1$\,Hz
        with a 1\,s timeout. Any \texttt{DEGRADED}/\texttt{DOWN} response,
        or any failure to respond within the timeout, causes a
        \texttt{INTERNAL\_SUBSYSTEM\_CRASH} (gateway) or
        \texttt{COMMUNICATION\_CHANNEL\_DOWN} (broker) failure event to be
        raised on \texttt{Av2DetectedFailures} -- which guarantees the
        ``\textsc{pms} detects any failures of internal communication
        subsystems within 5\,s'' response measure.
\end{enumerate}
Inside the \texttt{Av2CommunicationGateway} the responsibility is split
further: \texttt{Av2RequestReceiver} parses,
\texttt{Av2AuthenticationValidator} validates the device token,
\texttt{Av2RoutingManager} forwards either to the broker (sensor data) or
to the emergency path, \texttt{Av2FailureDetector} aggregates
intra-gateway failures and feeds the \texttt{Av2RetryCoordinator} that
re-attempts delivery with exponential back-off. The
\texttt{Av2GatewayRegistry} is the authoritative
\textsf{GatewayId}$\leftrightarrow$\textsf{PatientId}$\leftrightarrow$
expected-interval mapping and is updated whenever a
\texttt{P1IRiskEvents} message changes a patient's risk level.

\paragraph{(AD-Av2-3) Heartbeat-and-health detection within Av2 bounds.}
The QAS prescribes two distinct detection regimes:
\begin{itemize}
  \item \emph{Missing sensor data:} ``Detection time depends on the
        transmission rate configured on the patient gateway (and indirectly
        on the risk level), but does not exceed this with more than 5
        minutes.'' We implement this with a heartbeat tactic.
        \texttt{Av2RequestReceiver} calls
        \texttt{Av2HeartbeatUpdate.recordTransmission} for \emph{every}
        inbound transmission, regardless of authentication outcome, so
        connectivity is tracked independently of credential validity. The
        \texttt{Av2HeartbeatTracker} stores the last-seen timestamp and
        the configured expected interval per gateway, and the
        \texttt{Av2MissingDataDetector} (inside
        \texttt{Av2MonitoringService}) polls
        \texttt{Av2AvailabilityMonitoring} more frequently than the
        tightest expected interval, comparing
        $\textsf{now} - \textsf{lastHeartbeat}$ against
        \texttt{expectedInterval}\,+\,5\,min. The 5\,min slack is the
        explicit budget from the QAS.
  \item \emph{Internal subsystem failure:} ``The \textsc{pms} detects any
        failures of internal communication subsystems within 5 seconds
        after the failure.'' The \texttt{Av2InternalFailureDetector} polls
        the \texttt{Av2CommunicationGatewayHealthCheck} and
        \texttt{Av2MessageBrokerHealthCheck} endpoints. A polling
        frequency of $\geq 1$\,Hz with a 1\,s response timeout bounds
        total detection latency to under 5\,s; the polling frequency is
        therefore the sensitivity point for this measure (lowering it past
        1\,Hz threatens the deadline; raising it consumes CPU and network
        bandwidth that compete with sensor-data traffic).
\end{itemize}

A subtle point is that the \emph{expected interval} per gateway changes
when a patient's risk level changes (high-risk patients transmit more
often, so the missing-data threshold tightens accordingly). The
\texttt{Av2GatewayRegistry} subscribes to \texttt{P1IRiskEvents} and
propagates \texttt{updateExpectedInterval} into the
\texttt{Av2HeartbeatTracker}. This is the single ingress point through
which P1's risk-level updates affect Av2 detection sensitivity -- and one
of the cross-QAS sensitivity points discussed in
Chapter~\ref{ch:discussion}.

\paragraph{(AD-Av2-4) Risk-aware failure classification and notification.}
The QAS prescribes notification deadlines that depend on the risk level of
the affected patient: 2\,min (high), 5\,min (medium), 10\,min (low) and
that these be met in $\geq 90\,\%$ of cases. We separate three concerns:
\begin{itemize}
  \item \emph{Detection} produces a raw \texttt{FailureEvent} (one of
        \texttt{MISSING\_SENSOR\_DATA},
        \texttt{INTERNAL\_SUBSYSTEM\_CRASH},
        \texttt{COMMUNICATION\_CHANNEL\_DOWN}, \texttt{SLA\_VIOLATION});
  \item \emph{Classification} (\texttt{Av2FailureClassifier}) reads the
        affected patient's risk level via
        \texttt{Av2PatientStatusOverview} (provided by
        \texttt{PatientRecordRepository}) and stamps a
        \texttt{ClassifiedFailure} with an \texttt{EscalationLevel} and a
        notification deadline timestamp;
  \item \emph{Dispatch} (\texttt{Av2StakeholderNotifier} inside
        \texttt{Av2MonitoringService}) uses \texttt{Av2NotificationDispatch}
        on the dedicated \texttt{Av2StakeholderNotificationService} to fan
        out to the administrator (for subsystem/channel failures) and/or
        the patient's registered contact (for missing-data / gateway-down
        failures), with the contact phone number resolved through
        \texttt{Av2PatientContactLookup}.
\end{itemize}
Separating classification from dispatch (rather than collapsing them into
one ``alert builder'') is what permits the same dispatcher to be reused
for SLA-violation notifications without redoing the deadline arithmetic,
and is what allows the \texttt{Av2RecoveryTracker} to subscribe to the
same \texttt{Av2EscalationEvents} stream to start the recovery-coordination
clock.

\paragraph{(AD-Av2-5) Administrator-controlled recovery.}
Av2 distinguishes three resolution paths (failing subsystem, failing
channel, failing gateway). The architecture supports each:
\begin{itemize}
  \item \emph{Failing subsystem.} The administrator receives the
        notification through \texttt{Av2AdminPortal} and triggers
        \texttt{triggerRedeployment} or \texttt{triggerRollback} on the
        \texttt{Av2RedeploymentCoordinator} via
        \texttt{Av2RedeploymentControl}. The coordinator returns a
        \texttt{RedeploymentJob} handle and the operation is bounded to
        complete within 5\,min as required.
  \item \emph{Failing channel.} The administrator contacts the telecom
        operator out-of-band; the SLA with the operator stipulates support
        within 10\,min and resolution within the hour. The
        \texttt{Av2SLAMonitor} keeps the rolling availability / coverage /
        bandwidth statistics that justify SLA escalation.
  \item \emph{Failing patient gateway.} The
        \texttt{Av2StakeholderNotifier} dispatches a contact notification
        through \texttt{OutboundContactNotification} so the patient (or
        their contact) can resolve the issue (e.g.\ recharge the device).
\end{itemize}
On resolution, \texttt{Av2FailureClassifier} signals
\texttt{Av2RecoveryStatus} to the \texttt{Av2RecoveryTracker}, which
(a)~closes the failure record, (b)~updates the
\texttt{Av2CommunicationGapLog} with the gap duration, and (c)~initiates
\texttt{initiateRecovery} on \texttt{Av2RecoveryCoordination}. The
\texttt{Av2RecoverySyncService} then drains any data accumulated in the
broker during the outage and issues
\texttt{Av2GatewayResyncCommand.triggerResync} to each affected gateway,
which causes \texttt{Av2RetryAndSynchronizationService} on the gateway to
flush the \texttt{Av2LocalBufferingRepository} in chronological order.

\paragraph{(AD-Av2-6) SLA observability.}
The SLA stipulations (99.5\,\% availability, 80\,\% coverage, 64\,Kb/s
bandwidth) are externally negotiated; their architectural counterpart is
\emph{observability}, so that violations can be detected, evidenced, and
escalated. The \texttt{Av2SLAMonitor} reads aggregated metrics from
\texttt{Av2BandwidthMonitor} and the per-gateway heartbeat statistics
exposed by \texttt{Av2AvailabilityMonitoring}, and produces
\texttt{SLAReport} objects accessible to the \texttt{Av2AdminPortal} via
\texttt{Av2SLAReporting}. When any rolling threshold is breached the
monitor raises a \texttt{SLA\_VIOLATION} \texttt{FailureEvent}, which
feeds back into the same classification and notification pipeline.

\subsection*{Sequence diagrams}
The runtime behaviour of the four most quality-relevant flows is captured
by the sequence diagrams in the process view appendix: the
\emph{primary-path transmission with ack-and-clear}, the
\emph{degraded-mode transition with SMS escalation of a red emergency},
the \emph{missing-data detection / risk-classified notification path},
and the \emph{post-failure recovery / broker drain / gateway resync} flow.
All four diagrams go beyond the normal functional flow and explicitly
model the detection of exceptional situations, the system's response, and
the mechanisms enforcing the response measures (timeouts on the health
checks, retry deadlines on the routing manager, ack deadlines on the
acknowledgement service, the 5\,min slack on missing-data detection, and
the 2/5/10\,min deadline on the stakeholder notifier).

\subsection*{Considered alternatives}

\paragraph{Direct point-to-point gateway-to-ingestion connections.}
The simplest alternative is to let every patient gateway open a direct
connection to the \texttt{Av2DataIngestionService}, removing the
\texttt{Av2CommunicationGateway} and \texttt{Av2MessageBroker} from the
path. This was rejected for two reasons: (i)~it removes the single
observation point through which heartbeats can be tracked uniformly, which
would force the heartbeat tactic to be replicated inside the ingestion
service and cross-cut every storage write; (ii)~it couples gateway
availability to the availability of the persistence layer, so a transient
persistence outage would directly hit the gateways and undermine response 6
of P1. The gateway/broker pair pays a small per-message latency for
end-to-end decoupling and a single, uniform detection surface.

\paragraph{Persistent connection with TCP keep-alive instead of
heartbeats.}
TCP keep-alive could in principle serve as the liveness signal between
gateway and backend. We rejected it because keep-alive timeouts are set
per socket and per operating-system stack, are typically in the minutes
range, and cannot be calibrated per patient risk level. The QAS explicitly
ties detection latency to the transmission interval, which itself depends
on risk level. An application-level heartbeat that piggybacks on every
transmission (any inbound message counts as a heartbeat in
\texttt{Av2RequestReceiver}) is the only mechanism that meets the
$T{+}5$\,min measure across all risk tiers without artificially shortening
TCP keep-alive on every device.

\paragraph{Always-on backup channel for sensor data.}
Routing all sensor data via SMS would technically maximise availability,
but SMS bandwidth is well below 64\,Kb/s, costs scale linearly with
volume, and message ordering is not guaranteed. We restrict the backup
channel to emergency notifications, which are small, high-priority, and
rare, and rely on the \texttt{Av2LocalBufferingRepository} (24\,h
capacity) for the bulk path. This is consistent with the QAS resolution
responses, which only mandate the back-up channel for emergencies.

\paragraph{Auto-recovery of failed subsystems instead of administrator
control.}
A more aggressive alternative is to let the
\texttt{Av2RedeploymentCoordinator} redeploy autonomously upon any DOWN
status. We rejected this in favour of administrator-triggered recovery:
a misclassified failure (e.g.\ a transient \texttt{DEGRADED} caused by a
GC pause inside the broker) would cause unnecessary restarts, themselves
potentially destabilising the system at the worst possible moment. The
5\,min recovery budget is wide enough to permit human inspection of the
\texttt{Av2AdminPortal} before triggering the redeploy. The architecture
remains open to graduating to auto-recovery if the operational experience
proves the risk of false positives to be low.

\paragraph{Two-phase commit between gateway and backend instead of
ack-and-clear.}
A 2PC protocol would give strict-once semantics. The cost is a synchronous
round trip on every reading and elevated complexity in the
\texttt{Av2LocalBufferingRepository}. We accept at-least-once delivery
(with idempotent storage on the ingestion side keyed by gateway,
timestamp, and sequence) because (i)~sensor readings are immutable and
(ii)~duplicate detection at the storage layer is far cheaper than a 2PC
round trip on every reading -- and 2PC offers no advantage in the
catastrophic-failure scenario the QAS actually targets.


\section{P1: Data exchange with physicians (M)}
\label{sec:p1}

\subsection*{Key Decisions}
\begin{description}
  \item[(AD-P1-1)] \textbf{Thin physician edge gateway over
        command/query-split services.}
        The \texttt{P1PhysicianGateway} is a single external entry point
        with no business logic; behind it sit a write-side
        \texttt{P1PhysicianCommandService} (UC7/UC8/UC9), a read-side
        \texttt{P1PatientQueryService} (UC6) and a notification
        \texttt{P1NotificationDispatcher} (UC5).
  \item[(AD-P1-2)] \textbf{Risk-tiered priority queue on the read path.}
        \texttt{P1PatientQueryService} dispatches each
        \texttt{consultPatientStatus} call to a high/medium/low SLA tier
        derived from an in-memory \emph{tier index} that is kept current
        by subscribing to \texttt{P1IRiskEvents}. Unknown patients default
        to HIGH so that the SLA \emph{fails safe}.
  \item[(AD-P1-3)] \textbf{Read-through cache with high-risk pinning.}
        \texttt{P1PatientStatusCache} sits in front of the authoritative
        \texttt{P1PatientDataService} and pins high-risk patients in
        memory so that the 2\,s deadline is reachable under normal load.
  \item[(AD-P1-4)] \textbf{Event-driven notifications with priority buffer
        and durable log.}
        \texttt{P1NotificationDispatcher} subscribes to
        \texttt{P1IRiskEvents}, resolves recipients through the
        \texttt{P1RecipientRegistry}, persists red and yellow
        notifications in the \texttt{P1DurableNotificationLog}
        \emph{before} acknowledging the event, and enqueues into the
        \texttt{P1PriorityBuffer} which orders red\,$>$\,yellow\,$>$\,green
        and drops green first under overload.
  \item[(AD-P1-5)] \textbf{Correlation-tracked on-demand consultations
        (UC9).}
        \texttt{P1PhysicianCommandService} fetches fresh sensor data via
        \texttt{P1PatientGatewayCommander}, stores it with
        \texttt{triggerEstimation=false} to suppress the implicit
        scheduled job, then launches a \emph{priority} risk job carrying
        a \texttt{CorrelationId}. The result is pushed to the originating
        physician via \texttt{P1IConsultationResultPush}, no polling.
  \item[(AD-P1-6)] \textbf{Stale-tolerant EHR proxy.}
        \texttt{P1EHRProxyModule} fronts the external \texttt{HIS} with a
        cache invalidated on \textsc{pms}-initiated writes only. This is a
        documented sensitivity point with respect to externally-initiated
        HIS changes.
\end{description}

\subsection*{Employed tactics and patterns}
\textbf{Performance tactics:} \emph{prioritization} (risk-tiered queue,
priority buffer), \emph{introduce concurrency} (separate command/query
services), \emph{maintain multiple copies of computations} (status cache),
\emph{bound queue size} (priority buffer with drop-green-first policy),
\emph{event-driven notifications} (eliminate polling on UC9 results and
on status reads).
\textbf{Reliability tactics:} \emph{transactional log} (durable
notification log for red/yellow before ack), \emph{state resynchronization
on restart} (replay log into priority buffer).
\textbf{Patterns:} \emph{gateway} (physician gateway),
\emph{Command-Query Responsibility Segregation} between
\texttt{P1PhysicianCommandService} and \texttt{P1PatientQueryService},
\emph{publish/subscribe} on \texttt{P1IRiskEvents}, \emph{cache-aside /
read-through} for status and EHR, \emph{correlation identifier} for
asynchronous request/response in UC9.

\subsection*{Architectural rationale}

\paragraph{(AD-P1-1) Thin physician edge gateway over command/query-split
services.}
The \texttt{P1PhysicianGateway} exposes \texttt{P1IPhysicianAPI} and
dispatches each call to one of the three behind-the-gateway services
without performing any business logic. This is the standard \emph{gateway}
pattern, motivated here by three concerns:
\begin{enumerate}
  \item \emph{Independent scalability.} UC6 reads, UC7/8/9 writes, and
        UC5 notifications have different rate envelopes (25/min, 20/min
        and 100/min respectively) and different concurrency
        profiles. Co-locating them behind one service would force a
        single thread/connection budget to absorb the worst case of all
        three.
  \item \emph{Per-use-case timing.} The 2/5/10\,s deadlines (UC6), the
        3\,min initiation deadline (UC9 risk estimation), the 1\,min EHR
        forwarding deadline (UC8) and the 10\,s / 30\,s / 1\,min red /
        yellow / green dispatch deadlines (UC5) are all distinct;
        isolating them into separate services keeps each service's hot
        path simple enough to reason about.
  \item \emph{Failure containment.} Should the notification path be
        saturated, that congestion must not propagate to UC6 reads or UC9
        writes. The separate services share only the
        \texttt{P1IRiskEvents} pub/sub channel; back-pressure on one
        consumer (e.g.\ the notification dispatcher under a burst) does
        not block the others.
\end{enumerate}

\paragraph{(AD-P1-2) Risk-tiered priority queue on the read path.}
The QAS imposes \emph{three different} latency SLAs on UC6 depending on
the patient's current risk level (2\,s high, 5\,s medium, 10\,s low) and
a sustained throughput target of 25 requests/min. A single FIFO queue
cannot meet all three at once: under any burst exceeding 25/min, all
requests would be served strictly first-come-first-served and high-risk
reads would miss their 2\,s deadline behind a backlog of low-risk reads.

We therefore implement a three-tiered priority queue inside
\texttt{P1PatientQueryService}, fed from an in-memory \emph{tier index}
that maps \texttt{PatientId} $\rightarrow$ \texttt{RiskTier}. The tier
index is kept current by subscribing to \texttt{P1IRiskEvents}: every
status change pushed by \texttt{P1PatientStatusModule} (see AD-P1-4)
updates the index in $O(1)$. The hot path of a UC6 read is then a single
hash lookup followed by an enqueue onto the correct priority tier.
$O(1)$ on the hot path is the sensitivity point of this decision.
Unknown patients default to the \emph{HIGH} tier so that the SLA fails
safe when the index is cold (e.g.\ right after a restart).

\paragraph{(AD-P1-3) Read-through cache with high-risk pinning.}
\texttt{P1PatientStatusCache} sits between \texttt{P1PatientQueryService}
and \texttt{P1PatientStatusModule} (the authoritative status owner
inside \texttt{P1PatientDataService}), exposing
\texttt{P1IPatientStatusRead}. The cache uses a read-through strategy:
on miss, it consults \texttt{OtherDataMgmt} (provided by
\texttt{P1PatientStatusModule}) and populates itself. High-risk patients
are \emph{pinned} (never evicted) so the 2\,s deadline does not race a
cache fill on a cold key. Pinning is essentially a memory-for-time
trade-off. Freshness is bounded by a short TTL on cached entries: the
authoritative status lives in \texttt{P1PatientStatusModule} and any
write through \texttt{setEstimatedPatientStatus} is visible to the next
read after the TTL expires. We accept a brief staleness window on
status reads because the event-driven path that physicians actually
care about for timeliness, the UC5 notification fan-out, runs off
the \texttt{P1IRiskEvents} channel directly and does not depend on the
cache.

\paragraph{(AD-P1-4) Event-driven notifications with priority buffer and
durable log.}
UC5 imposes the most timing-sensitive response measures of P1:
red $\leq 10$\,s, yellow $\leq 30$\,s, green $\leq 1$\,min, sustained
$\geq 100$\,notifications/min, and \emph{no red or yellow notification
may be lost} under overload. The architecture combines three mechanisms
to meet all of them simultaneously:
\begin{itemize}
  \item \emph{Event-driven on the producer side.}
        \texttt{P1NotificationDispatcher} provides
        \texttt{P1IRiskEventListener} and subscribes to
        \texttt{P1IRiskEvents}. The single source of truth for ``a status
        changed'' is
        \texttt{P1PatientStatusModule.setEstimatedPatientStatus}, which
        fires \texttt{P1IRiskEvents} exactly once per real change. The
        dispatcher therefore reacts to real status changes instead of
        polling \texttt{P1IPatientQuery} on a timer, eliminating the
        polling load that would otherwise compete with the UC6 read
        path for capacity. Physicians retrieve queued notifications from
        the dispatcher's per-physician inbox
        (\texttt{P1INotificationInbox.getPendingNotifications} /
        \texttt{acknowledgeNotification}).
  \item \emph{Write-ahead durable log.} \texttt{P1RiskEventSubscriber}
        writes red and yellow notifications into the
        \texttt{P1DurableNotificationLog} \emph{before} the in-memory
        \texttt{P1PriorityBuffer} accepts them, so a crash between
        receipt and dispatch cannot lose a red/yellow notification. On
        restart, the log is replayed back into the buffer
        (\texttt{recoverPending}). Green notifications skip the log
        because their loss is tolerable under overload.
  \item \emph{Bounded priority buffer with drop policy.}
        \texttt{P1PriorityBuffer} is an in-memory priority queue ordering
        red\,$>$\,yellow\,$>$\,green and enforcing the 100/min throughput
        envelope. When the envelope is exceeded the buffer drops greens
        first; reds and yellows are protected by the durable log, so even
        if the buffer is briefly oversubscribed those entries are not
        lost (they remain in the log and are re-enqueued on the next
        admission cycle). This precisely satisfies ``it prioritizes red
        notifications over yellow and green ones, and yellow
        notifications over green ones. No red or yellow notifications may
        be lost.''
  \item \emph{Severity-tied dispatch deadlines.} The dispatcher reads
        each notification's severity and times out its delivery against
        the red/yellow/green deadlines (10/30/60\,s). If the primary
        delivery channel is unavailable, the
        \texttt{P1RiskEventSubscriber} can escalate red notifications
        onto the \texttt{Av2EmergencyDispatch} backup path. This is the
        explicit reuse point between Av2 and P1.
\end{itemize}

\paragraph{(AD-P1-5) Correlation-tracked on-demand consultations.}
UC9 (\emph{perform on-demand consultation}) is the most architecturally
demanding use case in P1: a single physician request must (i) fetch fresh
sensor data from the patient gateway, (ii) initiate a risk estimation
that is \emph{prioritised} over the routine scheduled jobs governed by
P2, (iii) return the result to the originating physician within 1\,min of
completion, and (iv) initiate the estimation itself within 3\,min of the
request, regardless of the patient's risk level. The flow is orchestrated
inside \texttt{P1OnDemandConsultationHandler}:
\begin{enumerate}
  \item Mint a fresh \texttt{CorrelationId} via
        \texttt{P1CorrelationTracker}.
  \item Pull fresh sensor data from the gateway via
        \texttt{P1IOnDemandSensorFetch} on
        \texttt{P1PatientGatewayCommander} (synchronous with a bounded
        timeout sized inside the 3\,min budget).
  \item Persist the readings through
        \texttt{SensorDataMgmt.addSensorData} with
        \texttt{triggerEstimation=false}, which is the key extension to
        the original \texttt{SensorDataMgmt} contract: without this flag,
        every UC9 call would trigger an extra \emph{implicit} scheduled
        risk-estimation job, doubling scheduler load and competing with
        P2's overload-mode budget for no functional benefit.
  \item Launch the risk estimation explicitly via
        \texttt{LaunchRiskEstimation} with \texttt{priority=HIGH} and the
        minted \texttt{CorrelationId}. The scheduler keeps the same
        FIFO/EDF policy from AD-2 of P2 but lets \texttt{HIGH} jump ahead
        of routine jobs, satisfying ``risk assessments triggered in an
        on-demand consultation are prioritized over regular risk
        estimations.''
  \item When the matching \texttt{RiskEvent} arrives on
        \texttt{P1IRiskEventListener},
        \texttt{P1CorrelationTracker.consumeCorrelation} matches it and
        the result is pushed back to the originating physician via
        \texttt{P1IConsultationResultPush} on
        \texttt{P1PhysicianGateway}.
\end{enumerate}
The end-to-end timing argument is: (a)~gateway fetch is bounded by the
commander timeout (well under 3\,min), (b)~scheduler initiation of a
priority job is constant-time once enqueued at the head, and (c)~the push
back to the physician uses an already-open channel (no extra connection
setup). The 1\,min ``result delivery'' deadline is therefore strictly
bounded by the propagation delay from \texttt{P1PatientStatusModule} to
\texttt{P1OnDemandConsultationHandler} via \texttt{P1IRiskEvents}, which
is sub-second under normal load.

\paragraph{(AD-P1-6) Stale-tolerant EHR proxy.}
\texttt{P1EHRProxyModule} fronts the external \texttt{HIS} via
\texttt{P1IHISAccess} and caches read results. UC8 requires that
``updates to a patient's EHR are forwarded within 1\,minute'': because
the proxy invalidates its cache on every \textsc{pms}-initiated write
(driven by \texttt{P1RiskLevelHandler}), the 1\,min deadline is bounded
by the round-trip to the HIS \texttt{savePatient}/
\texttt{saveRiskAssessment} operation. The proxy explicitly does
\emph{not} observe external updates to the HIS; this is documented as a
sensitivity point of P1 (Chapter~\ref{ch:discussion}) and is acceptable
in normal mode because the requirement scopes ``forwarding'' to
PMS-originated updates.

\subsection*{Sequence diagrams}
The process view contains a sequence diagram for each major P1 response:
the tier-classified UC6 read with cache hit/miss, the UC5 fan-out of a
red notification under and above the 100/min threshold, and the UC9
end-to-end correlation-tracked flow.

\subsection*{Considered alternatives}

\paragraph{Producer-side polling for status changes.}
An alternative architecture in which \texttt{P1NotificationDispatcher}
periodically polls \texttt{P1IPatientQuery} for every subscribed patient
to discover new status changes would satisfy the functional requirement
of UC5. We rejected it because the poll frequency required to meet the
10\,s red deadline (every $\leq 5$\,s per patient) multiplied by an
arbitrary number of subscribed patients makes the read load unbounded
in the worst case, directly conflicting with response 6 of P1 (no
impact on the risk-estimation pipeline). Event-driven \texttt{P1IRiskEvents}
subscription fires $O(1)$ per real status change, not per unit of time,
and decouples discovery latency from subscriber count. The physician
inbox itself remains pull-based on the consumer side
(\texttt{getPendingNotifications}); the alternative we reject here is
polling on the \emph{producer-internal} status-change discovery, not
the inbox surface.

\paragraph{Single FIFO queue for all UC6 reads.}
A single FIFO queue would deliver the 25\,reads/min throughput but cannot
honour the per-risk-tier deadlines. We discussed an EDF (earliest
deadline first) scheduler instead of three tiered queues; it gives the
same mathematical guarantee but requires the deadline computation on
every enqueue, putting additional pressure on the very hot path the SLA
is trying to protect. The tiered design pays a small memory cost (three
buffers) for an $O(1)$ classify-and-enqueue.

\paragraph{Single write-through cache for both status and EHR.}
Combining \texttt{P1PatientStatusCache} and the EHR cache inside
\texttt{P1EHRProxyModule} into one structure would simplify the component
graph but conflate two different consistency models: status is strongly
consistent (event-driven invalidation), EHR is stale-tolerant on reads.
Conflating them forces the cache to use whichever model is stricter,
which would either make EHR reads unnecessarily expensive or status
reads unnecessarily stale.

\paragraph{Synchronous UC9 (long-polling on the physician channel).}
Holding the physician's request open for the duration of the
risk-estimation job is the simplest implementation of UC9 but it (i)~ties
up an inbound connection for up to 3\,min, (ii)~hides the asynchronous
nature of the pipeline from the physician's client and prevents the
client from listening to other notifications in the meantime, and
(iii)~requires the \texttt{P1PhysicianGateway} to be stateful for the
duration. The correlation-id design keeps the gateway stateless and the
physician's channel free for concurrent UC5 notifications, a hard
requirement under the P1 throughput envelope.

\paragraph{Allowing UC9 to use the standard scheduled path.}
We considered routing UC9 through the same scheduler entry that handles
new sensor arrivals (the default of \texttt{LaunchRiskEstimation}). This
was rejected because UC9 explicitly requires that its job be
\emph{prioritised} over routine jobs ``independent of the patient's risk
level'', which the standard FIFO/EDF policy does not provide. Extending
\texttt{LaunchRiskEstimation} with an explicit \texttt{Priority}
parameter is the smallest possible change to the original P2 contract
that honours this requirement (see Chapter~\ref{ch:discussion}).


% ============================================================
% Part 3: Other (supporting) decisions
% ============================================================
\chapter{Other decisions}
\label{ch:other-decisions}

The decisions documented in Chapter~\ref{ch:qas-decisions} are
\emph{quality-driven}. Around them the extension also takes a small
number of \emph{supporting} functional decisions whose primary role is to
make the quality responses meaningful. We list them here for traceability;
their detailed responsibilities are documented in the element catalog
generated by the SAPlugin.

\paragraph{(OD-1) Replacement of \texttt{OtherFunctionality} by concrete
patient-data services.}
The initial architecture absorbed every responsibility outside the
risk-estimation pipeline into the placeholder component
\texttt{OtherFunctionality}. This placeholder is replaced by
\texttt{P1PatientDataService} (sensor data, patient status, HIS EHR proxy
via \texttt{P1HISAdapter}), the \texttt{PatientRecordRepository}
(authoritative patient identity and risk level),
\texttt{PatientManagementService} and \texttt{UserAuthenticationService},
and the dedicated \texttt{Av2SensorDataRepository} for raw sensor
packages. The three interfaces previously required from
\texttt{OtherFunctionality} (\texttt{OtherDataMgmt},
\texttt{PatientRecordMgmt}, \texttt{SensorDataMgmt}) are now provided by
\texttt{P1PatientDataService}, with unchanged operation contracts where
possible. The single change to the \texttt{SensorDataMgmt} contract is
the addition of an explicit \texttt{triggerEstimation} boolean on
\texttt{addSensorData}, motivated by AD-P1-5; UC4 callers (the default)
pass \texttt{true} to preserve the original behaviour, UC9 callers pass
\texttt{false} to suppress the implicit scheduled job in favour of an
explicit priority launch.

\paragraph{(OD-2) Extension of the \texttt{LaunchRiskEstimation} contract.}
The original \texttt{LaunchRiskEstimation} carried \texttt{patientId},
\texttt{newSensorData} and \texttt{receivedAt}. It is extended in three
ways: (i)~\texttt{Priority} is added so UC9 jobs can jump ahead of
routine ones (AD-P1-5); (ii)~\texttt{CorrelationId} is added (optional)
so the originating UC9 request can be matched to the resulting event;
and (iii)~\texttt{newSensorData} and \texttt{receivedAt} become
\emph{optional}, so a UC7 reconfigure can launch a re-assessment
against the patient's most recent stored sensor package without having
to re-supply it. The scheduling logic from AD-2 of P2 is unchanged:
FIFO in normal mode, EDF in overload mode, with the additional rule
that \texttt{Priority = HIGH} jumps ahead of routine jobs. The
\texttt{CorrelationId} parameter is opaque to the scheduler and is
propagated unchanged through to the \texttt{RiskEvent} payload, where
it allows subscribers to match the result back to the originating
request. All three changes are strictly additive at the interface
level: callers that pass \texttt{NORMAL}/\texttt{null} for the new
fields and supply \texttt{newSensorData}/\texttt{receivedAt} preserve
the original semantics.

\paragraph{(OD-3) Introduction of \texttt{P1IRiskEvents} as the single
in-process event bus.}
\texttt{P1IRiskEvents} is emitted by \texttt{P1PatientStatusModule}
whenever \texttt{setEstimatedPatientStatus} \emph{actually changes} the
stored status. Subscribers include \texttt{P1NotificationDispatcher} (UC5
fan-out via \texttt{P1RiskEventSubscriber}),
\texttt{P1PatientQueryService} (tier-index refresh), and
\texttt{P1OnDemandConsultationHandler} (UC9 result delivery).
Centralising risk-state propagation onto one event topic prevents an
$N{\times}M$ matrix of point-to-point notifications and is the
architectural pre-condition for the tier index of AD-P1-2 as well as
the recipient fan-out and UC9 correlation match of AD-P1-4 and
AD-P1-5.

\paragraph{(OD-4) Three deployment profiles.}
The deployment view retains the \emph{primary} (logical), \emph{pilot}
(200 patients, 5 clinical models) and \emph{dev} (20 patients,
3 clinical models) profiles from the initial architecture and extends
them with the new Av2/P1 nodes (\texttt{GatewayNode},
\texttt{BackendGatewayNode}, \texttt{MonitoringNode}, \texttt{AdminNode},
\texttt{PatientRecordsNode}, \texttt{PatientDataNode},
\texttt{PhysicianAccessNode}, \texttt{SMSGatewayProvider},
\texttt{TelecomNetwork}, \texttt{HISServer}). The pilot profile
co-locates the monitoring tier on its own node so that a failure in the
ingress chain does not also take down the detection mechanism; the dev
profile collapses everything onto one back-end host to keep local
testing tractable. The placement of the
\texttt{Av2RedeploymentCoordinator} on a separate \texttt{AdminNode} is
deliberate: the recovery control plane must remain reachable when the
data plane is degraded.


% ============================================================
% Part 4: Discussion -- trade-offs and sensitivity points
% ============================================================
\chapter{Discussion}
\label{ch:discussion}

This chapter discusses the impact of the extension on the initial
architecture. It is organised around the three required analyses of
Section~1.2 of the assignment: (i)~impacted decisions from the initial
architecture, (ii)~the relationship between the extension and each
impacted decision (reinforce / weaken / revise), and (iii)~the
acceptability of the trade-offs in the \textsc{pms} context.

\section{Impacted decisions from the initial architecture}

The initial architecture documented four QAS decisions:
\begin{itemize}
  \item P2-AD1: \emph{Parallel and concurrent risk estimation} across
        multiple \texttt{RiskEstimationProcessor} instances and multiple
        nodes;
  \item P2-AD2: \emph{Centralized scheduling with FIFO/EDF overload mode}
        in \texttt{RiskEstimationScheduler};
  \item M2-AD1: \emph{Run-time deployment of clinical models} via an
        append-only \texttt{ClinicalModelDB};
  \item M2-AD2: \emph{Local caching with explicit invalidation} of
        clinical models and their configurations in
        \texttt{ClinicalModelCache}.
\end{itemize}
The extension also implicitly relies on the catch-all
\texttt{OtherFunctionality} placeholder being well-behaved, which is the
fifth impacted assumption.

\section{Relationship to existing decisions}

\paragraph{P2-AD1 (parallel risk estimation) is \emph{reinforced} but
more tightly constrained.}
The parallel processor pool is unchanged in its responsibilities. What
changes is the source mix on the scheduler's input queue: in addition to
the routine arrivals previously delivered through the
\texttt{OtherFunctionality} placeholder, the queue now receives
\texttt{Priority=HIGH} jobs from \texttt{P1OnDemandConsultationHandler}
(UC9) and \texttt{Priority=NORMAL} jobs from
\texttt{P1ConfigurationHandler} (UC7) and \texttt{P1SensorIngestModule}
(UC4 default). The processor pool is therefore \emph{reinforced} (its
purpose is now more concrete) but \emph{more tightly constrained}
(the worst-case arrival pattern is wider because UC9 priority jobs can
preempt routine ones, increasing the variability of effective per-job
latency for low-priority jobs).

\paragraph{P2-AD2 (centralized scheduling) is \emph{revised additively}.}
The \texttt{LaunchRiskEstimation} contract gains the \texttt{Priority}
and \texttt{CorrelationId} parameters (OD-2). The underlying FIFO/EDF
logic is unchanged; only the additional priority class is added. We
deliberately did \emph{not} change the original overload threshold (20
risk estimations/min) or the deadline table (2/5/8 min) for normal and
overload mode. The \texttt{HIGH} priority class is orthogonal to the
overload mode: a UC9 job with \texttt{HIGH} priority jumps the queue
regardless of mode, and its initiation is bounded by the 3\,min UC9
measure. This is a \emph{revision} because it changes the observable
behaviour of the queue for routine jobs (a routine job may now be passed
by a priority job), but the revision is strictly additive at the
interface level and at the policy level: a system that issues only
\texttt{NORMAL}-priority calls is indistinguishable from the original
system.

\paragraph{M2-AD1 (run-time model deployment) is \emph{preserved
unchanged}.}
The \texttt{ClinicalModelDB}, \texttt{ClinicalJobCreator},
\texttt{RiskEstimationProcessor}, \texttt{MLModelManager}, and the entire
MLaaS sync / update-processing pipeline (\texttt{MLModelUpdateProcessor},
\texttt{ModelUpdater}, etc.) are untouched. The extension does not
introduce any new path through the clinical model machinery; it only
adds a new \emph{caller} of the existing \texttt{ClinicalModelCacheMgmt}
interface, namely \texttt{P1ConfigurationHandler} (UC7).

\paragraph{M2-AD2 (clinical-model caching) is \emph{extended in scope}.}
The original cache invalidation was triggered by the
\texttt{OtherFunctionality} placeholder. After the extension the trigger
moves to the explicit \texttt{P1ConfigurationHandler}, which calls
\texttt{invalidateCacheEntries} after a successful UC7 write via
\texttt{P1ClinicalConfigMgmt}. The cache semantics are preserved
(snapshot stability for in-flight jobs, fresh data for newly launched
jobs). The cache is also referenced by the new
\texttt{P1PatientStatusModule} only indirectly (through the status path,
not the clinical-model path), so M2-AD2's reasoning about in-flight job
correctness is preserved verbatim.

\paragraph{The \texttt{OtherFunctionality} assumption is
\emph{abandoned}.}
The initial architecture explicitly documented the placeholder component
as ``other functionality that still has to be worked out''. The extension
fully discharges this placeholder by introducing
\texttt{P1PatientDataService}, \texttt{PatientRecordRepository},
\texttt{PatientManagementService}, \texttt{UserAuthenticationService},
and \texttt{Av2SensorDataRepository} (OD-1). All three previously
required-from-placeholder interfaces (\texttt{OtherDataMgmt},
\texttt{PatientRecordMgmt}, \texttt{SensorDataMgmt}) are now provided by
concrete components and the corresponding diagram boxes are no longer
greyed out.

\section{Sensitivity points introduced or affected}

The following sensitivity points were introduced or aggravated by the
extension. ``Sensitivity point'' is used in the ATAM sense: an
architectural decision (or parameter) for which a small change has a
disproportionate impact on one or more quality attributes.

\begin{itemize}
  \item \textbf{Polling frequency of
        \texttt{Av2InternalFailureDetector}.}
        Detection of internal subsystem failures is sensitive to the
        polling frequency. At $\geq 1$\,Hz the 5\,s response measure is
        achievable; below 0.5\,Hz it is at risk. The poll cost competes
        with other CPU and network budget on \texttt{MonitoringNode}.
  \item \textbf{Configured transmission interval per gateway.}
        The Av2 missing-data detection deadline ($T{+}5$\,min) is
        \emph{relative} to the configured interval. Shortening the
        interval for high-risk patients tightens the detection window
        but also increases load on \texttt{Av2CommunicationGateway},
        \texttt{Av2MessageBroker} and the downstream ingestion pipeline.
        The interval is therefore a sensitivity point \emph{shared} with
        P2 (it bounds the arrival rate the scheduler must absorb).
  \item \textbf{Cache hit rate of \texttt{P1PatientStatusCache}.}
        The 2\,s read deadline for high-risk patients is sensitive to
        the cache hit rate. Pinning high-risk patients defends the hot
        case; but if the high-risk population grows beyond cache
        capacity, the hit rate degrades and the 2\,s deadline becomes
        contended. The cache size is a tunable sensitivity point.
  \item \textbf{Size of the priority buffer inside
        \texttt{P1NotificationDispatcher}.}
        Too small: green notifications are dropped earlier than
        necessary. Too large: memory consumption grows under sustained
        overload and red/yellow dispatch latency degrades because the
        buffer takes longer to drain.
  \item \textbf{Durability semantics of the
        \texttt{P1DurableNotificationLog}.}
        Whether a write is \texttt{fsync}-ed before the event is acked
        determines whether a red notification can be lost on a crash.
        The ``no red or yellow may be lost'' response measure is
        sensitive to this single configuration decision.
  \item \textbf{SLA threshold definitions in \texttt{Av2SLAMonitor}.}
        The 99.5\,\% / 80\,\% / 64\,Kb/s figures are external SLA
        parameters; the rolling window over which they are evaluated is
        an internal architectural choice. Too short a window produces
        false positives during normal transient outages; too long a
        window masks real violations.
  \item \textbf{Single-topic coupling: \texttt{P1IRiskEvents}.}
        A single in-process event topic carries every status change.
        Three P1 consumers subscribe to it (notification fan-out,
        query-service tier index, UC9 correlation match). Any added
        latency on this topic propagates to all of them, so the topic
        is a single point of sensitivity for P1 throughput SLAs.
\end{itemize}

\section{Trade-off points}

The following decisions are \emph{trade-off points} in the ATAM sense:
they improve one quality at the explicit cost of another. In the
headings below, $\uparrow$ marks a quality that the decision
\emph{improves} and $\downarrow$ marks a quality that the decision
\emph{degrades}.

\paragraph{Broker-decoupled ingress (AD-Av2-2): availability $\uparrow$
vs.\ end-to-end latency $\downarrow$.}
Inserting the \texttt{Av2MessageBroker} between the gateway and the
ingestion service adds a hop on every reading. Average end-to-end
latency from sensor reading to persisted record increases by the broker
write latency. In exchange we get failure isolation between ingress and
persistence (the broker absorbs transient persistence outages) and
durability for the ack-and-clear protocol with the gateway. The
trade-off is acceptable because Av2 is a high-priority scenario and
end-to-end ingestion latency is not a P2 deadline (the scheduler
operates on already-stored data; the broker write is amortised over the
much larger estimation time).

\paragraph{Centralised \texttt{Av2MonitoringService} (AD-Av2-3, AD-Av2-4):
predictable detection $\uparrow$ vs.\ single point of operational risk
$\downarrow$.}
Centralising failure detection and classification simplifies global
deadline enforcement and avoids inconsistent decisions across
distributed detectors. The cost is that the monitoring service is itself
critical infrastructure; its own failure must not be self-masking. We
mitigate this by deploying the monitoring service on a separate
\texttt{MonitoringNode}, by deploying the
\texttt{Av2RedeploymentCoordinator} on yet a separate
\texttt{AdminNode}, and by making
\texttt{Av2CommunicationGatewayHealthCheck} and
\texttt{Av2MessageBrokerHealthCheck} polled remotely. A monitoring
outage would itself be visible to the \texttt{Av2AdminPortal} via the
absence of \texttt{Av2SLAReporting} updates.

\paragraph{Tier-index in-memory cache (AD-P1-2): UC6 performance $\uparrow$
vs.\ consistency $\downarrow$ for a millisecond window.}
The tier index is updated reactively when \texttt{P1IRiskEvents} fires.
Between the underlying status change and the index update, a UC6 read
could be classified into the wrong tier. We accept this because the
window is sub-second under normal load and the tier classification only
affects the SLA tier (the read still returns the latest status from the
cache); the architectural ``fail-safe to HIGH'' default further bounds
the worst case.

\paragraph{Drop-green-first under overload (AD-P1-4): red/yellow
guarantee $\uparrow$ vs.\ green completeness $\downarrow$.}
At sustained $>100$ notifications/min, greens may be dropped before
reaching the physician. This is a deliberate trade-off \emph{prescribed}
by the QAS itself (``no red or yellow notifications may be lost'' is
the only completeness measure stated) and we adopt it verbatim.

\paragraph{Stale-tolerant EHR proxy (AD-P1-6): read performance $\uparrow$
vs.\ externally-initiated update visibility $\downarrow$.}
The cache is invalidated only on \textsc{pms}-initiated writes. External
updates made directly to the HIS are not observed until the cache entry
expires or is invalidated by an unrelated PMS write. This is acceptable
under the QAS, which only constrains \emph{forwarding} from \textsc{pms}
to the HIS, not the reverse path. If the \textsc{pms} ever needs to
observe externally-initiated updates in near-real-time, this trade-off
must be revisited (e.g.\ by adding a TTL or by subscribing to FHIR
subscriptions on the HIS).

\paragraph{Priority injection into the scheduler (OD-2, AD-P1-5): UC9
timing $\uparrow$ vs.\ predictability for routine jobs $\downarrow$.}
Allowing UC9 jobs to preempt routine ones improves the 3\,min UC9
initiation measure but increases the worst-case wait for routine jobs
under burst conditions. Under the P2 overload threshold (20/min) and
the P1 UC9 rate (20/min) the worst-case routine wait is bounded by the
combined arrival rate; in particular it does not invalidate the P2
10\,min hard deadline for any single risk-estimation job, which remains
the original safety floor.

\section{Acceptability of the trade-offs}

The trade-offs above are acceptable in the \textsc{pms} context for the
following reasons:
\begin{itemize}
  \item All cost-side qualities affected by the trade-offs are
        \emph{measured and bounded}: ingress latency by the broker SLA,
        tier-index staleness by the event-bus latency, green-drop rate
        by the priority-buffer policy, and routine-job preemption by
        the P2 hard deadline. Where a trade-off makes a guarantee less
        strict (greens may be dropped, tier index is briefly stale),
        the QAS \emph{itself} accepts that relaxation.
  \item Av2 is a \emph{High} priority scenario; P1 is \emph{Medium};
        and the original P2 (Medium) and M2 (High) guarantees are
        preserved. The priority ordering of trade-off resolutions
        therefore aligns with the stakeholder priorities: clinical
        safety (M2 correctness, P2 deadline) is never sacrificed;
        gateway availability (Av2) is never sacrificed; physician
        convenience properties (UC6 tier latency, UC5 green
        completeness) absorb the cost where one must be paid.
  \item Every trade-off is \emph{locally reversible}: changing the
        polling frequency, the cache size, the buffer size, the SLA
        window, or the durability semantics of the notification log
        does not require any structural change to the architecture.
        The sensitivity points are exposed as configuration parameters
        rather than baked into the topology.
  \item The single \emph{structural} commitment that is hard to reverse,
        the broker-decoupled ingress, is justified by the High
        priority of Av2 and is consistent with how the rest of the
        system already treats persistent state (compare M2's
        append-only \texttt{ClinicalModelDB}, which makes the same
        durability-over-latency choice for clinical models).
\end{itemize}

\section{Risks and limitations}

The architecture explicitly does not address the following residual
risks, which are out of scope for the selected QASs but are worth
recording for future iterations:
\begin{itemize}
  \item \emph{No autonomous recovery of failed subsystems.} Recovery
        requires administrator action via \texttt{Av2AdminPortal}. If
        the administrator is unreachable, the 5\,min recovery measure
        cannot be met. The QAS accepts this; a future extension could
        add a watchdog autonomous rollback gated by a cooling-off
        period.
  \item \emph{No security QAS coverage.} The
        \texttt{Av2AuthenticationValidator} (gateway-device
        authentication) and \texttt{UserAuthenticationService}
        (human-user authentication) are present in the architecture but
        only as artefacts of Av2 and the basic functional flow; the
        report does not derive their design from a security QAS.
  \item \emph{Single \texttt{Av2MessageBroker} instance.} The
        architecture treats the broker as a logical durable queue; the
        deployment view shows a single instance per environment.
        Should broker availability prove to be a bottleneck, the
        architecture can be extended to a replicated broker without
        changes to the surrounding components; this is left as an
        operational tuning concern rather than an architectural one.
  \item \emph{No cross-version dual-writes for clinical models.}
        M2-AD1 already documents the snapshot-stability argument for
        in-flight jobs; the extension does not change it.
\end{itemize}

\bigskip
The final architecture --- diagrams (system context, client-server,
decomposition, deployment, process view) and the full element catalog
--- is provided in the appendices, generated from the Visual Paradigm
project through the SAPlugin.
